\documentclass[12pt]{article}
\usepackage{amsmath, amssymb}
\usepackage{geometry}
\geometry{margin=1in}

\begin{document}

%%%%%%%%%%%%%%%%%%%%%%%%%%%%%%%%%%%%%%%%%%%%%%%%%%%%%%
\begin{center}
    \Large \textbf{Math 521 -- Practice Problems (Countability)}\\
    \normalsize Problem Set
\end{center}

\newpage
%%%%%%%%%%%%%%%%%%%%%%%%%%%%%%%%%%%%%%%%%%%%%%%%%%%%%%

\section*{Problem 1}

\textbf{(a)} State the definition of: the set $A$ is countable.

\vspace{1cm}

\textbf{(b)} Prove that the set
\[
B = \{(x,y) \in \mathbb{R}^2 : x \in \mathbb{N} \text{ and } y \in \{0,1,2\}\}
\]
satisfies the definition in part (a).

\vfill
\newpage

%%%%%%%%%%%%%%%%%%%%%%%%%%%%%%%%%%%%%%%%%%%%%%%%%%%%%%

\section*{Problem 2}

\textbf{(a)} State the definition of: the set $A$ is countable.

\vspace{1cm}

\textbf{(b)} Prove that the set
\[
C = \{(x,y) \in \mathbb{R}^2 : x \in 2\mathbb{N} \text{ and } y \in \{0,1\}\}
\]
is countable.

(Here $2\mathbb{N} = \{2n : n \in \mathbb{N}\}$.)

\vfill
\newpage

%%%%%%%%%%%%%%%%%%%%%%%%%%%%%%%%%%%%%%%%%%%%%%%%%%%%%%

\section*{Problem 3}

\textbf{(a)} State the definition of: the set $A$ is countable.

\vspace{1cm}

\textbf{(b)} Prove that the set
\[
D = \{(x,y) \in \mathbb{R}^2 : x \in \mathbb{N} \text{ and } y \in \mathbb{Z}, \ |y| \le 3\}
\]
is countable.

\vfill

%%%%%%%%%%%%%%%%%%%%%%%%%%%%%%%%%%%%%%%%%%%%%%%%%%%%%%
\newpage

\section*{Problem 4 (Field Practice)}

Let $(F,+,\cdot)$ be a field.

Prove that if $a \in F$ and $a \neq 0$, then $(-a)^{-1} = -(a^{-1})$.

Make sure each step follows from a property of fields stated in lecture.

\vfill

%%%%%%%%%%%%%%%%%%%%%%%%%%%%%%%%%%%%%%%%%%%%%%%%%%%%%%
\newpage

\section*{Problem 5 (Field Practice)}

Let $(F,+,\cdot)$ be a field.

Prove that if $a,b \in F$ and $ab = 0$, then $a = 0$ or $b = 0$.

(You may use the existence of multiplicative inverses for nonzero elements.)

\vfill

%%%%%%%%%%%%%%%%%%%%%%%%%%%%%%%%%%%%%%%%%%%%%%%%%%%%%%
\newpage

\section*{Problem 6 (Field Practice)}

Let $(F,+,\cdot)$ be a field.

Prove that if $a,b,c \in F$ and $a \neq 0$ and $ab = ac$, then $b = c$.

Make sure your argument clearly uses properties of fields.

\vfill

%%%%%%%%%%%%%%%%%%%%%%%%%%%%%%%%%%%%%%%%%%%%%%%%%%%%%%
\newpage

\section*{Problem 7 (Field Practice)}

Let $(F,+,\cdot)$ be a field.

Prove that the additive identity in $F$ is unique.

That is, if $e \in F$ satisfies $a + e = a$ for all $a \in F$, then $e = 0$.

\vfill

%%%%%%%%%%%%%%%%%%%%%%%%%%%%%%%%%%%%%%%%%%%%%%%%%%%%%%
\newpage

\section*{Problem 8 (Supremum Practice)}

Let $A, B \subset \mathbb{R}$ be nonempty and bounded above sets. Define
\[
S = \{a + b : a \in A,\ b \in B\}.
\]

Prove that $S$ is bounded above.

\vfill

%%%%%%%%%%%%%%%%%%%%%%%%%%%%%%%%%%%%%%%%%%%%%%%%%%%%%%
\newpage

\section*{Problem 9 (Supremum Practice)}

Let $A \subset \mathbb{R}$ be nonempty and bounded above, and let $c \in \mathbb{R}$. Define
\[
cA = \{ca : a \in A\}.
\]

Prove that if $c > 0$, then
\[
\sup(cA) = c \sup A.
\]

\vfill

%%%%%%%%%%%%%%%%%%%%%%%%%%%%%%%%%%%%%%%%%%%%%%%%%%%%%%
\newpage

\section*{Problem 10 (Infimum Practice)}

Let $A \subset \mathbb{R}$ be nonempty and bounded below, and let $c \in \mathbb{R}$ with $c > 0$. Define
\[
cA = \{ca : a \in A\}.
\]

Prove that
\[
\inf(cA) = c \inf A.
\]

\vfill

%%%%%%%%%%%%%%%%%%%%%%%%%%%%%%%%%%%%%%%%%%%%%%%%%%%%%%
\newpage

\section*{Problem 11 (Supremum / Infimum Practice)}

Let $A, B \subset \mathbb{R}$ be nonempty and bounded below sets. Define
\[
T = \{a + b : a \in A,\ b \in B\}.
\]

Prove that
\[
\inf T = \inf A + \inf B.
\]

\vfill

%%%%%%%%%%%%%%%%%%%%%%%%%%%%%%%%%%%%%%%%%%%%%%%%%%%%%%
\newpage

\section*{Problem 12 (Sequence Practice)}

Let $\{a_n\}_{n\ge 1}$ be a sequence and $a \in \mathbb{R}$.

\textbf{(a)} Prove or disprove: If $\{a_n\}$ converges to $a$, then $\{a_n^2\}$ converges to $a^2$.

\vfill

%%%%%%%%%%%%%%%%%%%%%%%%%%%%%%%%%%%%%%%%%%%%%%%%%%%%%%
\newpage

\section*{Problem 13 (Sequence Practice)}

Let $\{a_n\}_{n\ge 1}$ be a sequence and $a \in \mathbb{R}$.

\textbf{(a)} Prove or disprove: If $\{a_n^2\}$ converges to $a^2$, then $\{a_n\}$ converges to $a$.

\vfill

%%%%%%%%%%%%%%%%%%%%%%%%%%%%%%%%%%%%%%%%%%%%%%%%%%%%%%
\newpage

\section*{Problem 14 (Sequence Practice)}

Let $\{a_n\}_{n\ge 1}$ and $\{b_n\}_{n\ge 1}$ be sequences.

\textbf{(a)} Prove or disprove: If $a_n \to a$ and $b_n \to b$, then $a_n b_n \to ab$.

(You may use limit laws proven in lecture, but justify each step.)

\vfill

%%%%%%%%%%%%%%%%%%%%%%%%%%%%%%%%%%%%%%%%%%%%%%%%%%%%%%
\newpage

\section*{Problem 15 (Sequence Practice)}

Let $\{a_n\}_{n\ge 1}$ be a sequence.

\textbf{(a)} Prove or disprove: If $a_n \ge 0$ for all $n$ and $a_n \to 0$, then $\sqrt{a_n} \to 0$.

\vfill

%%%%%%%%%%%%%%%%%%%%%%%%%%%%%%%%%%%%%%%%%%%%%%%%%%%%%%
\newpage

\section*{Problem 16 (Sequence Practice)}

Let $\{a_n\}_{n\ge 1}$ be a sequence.

\textbf{(a)} Prove or disprove: If $|a_n| \to 0$, then $a_n \to 0$.

\vfill

\end{document}